\chapter{Team Contribution}
\label{app:contrib}
\begin{table}[H]
\centering
\caption{Components owned by each team member}
\small
\begin{tabularx}{\linewidth}{@{}l l X@{}}
\toprule
\textbf{USN} & \textbf{Name} & \textbf{Components owned} \\ \midrule
\USNa & \NameA\ (Team Lead) & Orchestration and integration; command-line interface; LangGraph workflow and model gateway; snapshots, calculation profile and approval control; shared schemas and configuration; trial coordinator, host sandbox and evaluation tooling; local model server \\
\USNb & \NameB & Source acquisition and intake; PDF, CSV and JSON extraction; index construction; retrieval and applicability; regulatory corpus content and citations \\
\USNc & \NameC & Normalisation and equivalence; deterministic domain tools and exact reconciler; synthetic world generator, renderers and fault injection; development and held-out dossiers; development-set tuning of the prompt on a faster machine \\
\USNd & \NameD & Blob store and catalog; audit and transaction kernel; reports and verifier; review service; export and replay; evaluator and oracle isolation; human rubric \\
\bottomrule
\end{tabularx}
\end{table}

The split above follows the implementation plan of the repository. Every member is expected to explain the whole
system. The report, slides and demonstration are shared among the four members.

\pending{Percentage shares of contribution, if the course requires them, and any change in ownership since the
plan was written, need to be confirmed by the team. The tuning of the prompt on the faster machine and the
held-out run have not been completed.}

\chapter{Setup and Run Instructions}
\label{app:run}
The repository is at \url{https://github.com/R-V-Abhishek/ReCycleGuard}. The branch with this work is
\texttt{checkpoint-5}. No API keys are needed, because the model is local and the endpoint is loopback only.
A secret such as the API token must be created locally and never committed.

\begin{lstlisting}[language={},numbers=none,caption={Portable part (any operating system, no model)},label=lst:run1]
git clone https://github.com/R-V-Abhishek/ReCycleGuard.git && cd ReCycleGuard
git checkout checkpoint-5
uv sync --locked
uv run pytest -m "not slow and not model"     # fast suite; full suite: uv run pytest
uv run ruff check . && uv run ruff format --check .
\end{lstlisting}

\begin{lstlisting}[language={},numbers=none,caption={API and interface (Apple silicon for real-model runs)},label=lst:run2]
# 1. native model (separate project, Apple silicon only)
uv run --project tools/local-model python tools/local-model/serve.py
# 2. API: set a random token of at least 24 characters first
export RECYCLEGUARD_API_TOKEN=<your token>
uv run recycleguard serve --host 0.0.0.0 --port 8000 \
  --model-config .recycleguard/local-model.json --cors-origin http://localhost:5173
# 3. user interface (Node 22 or later)
cd ui && npm ci && npm run dev -- --host
\end{lstlisting}

\begin{lstlisting}[language={},numbers=none,caption={Evaluation on the development dossiers},label=lst:run3]
uv run scripts/run_trials.py docs/checkpoints/mg1/protocols/reliability.json <work>/reliability \
  --config .recycleguard/local-model.json --venv <work>/wheel-venv
uv run scripts/score_trials.py <work>/reliability <oracle-root> tests/fixtures/visible \
  <work>/reliability-scores
uv run scripts/make_tables.py <work>/reliability <work>/reliability-scores <work>/reliability-tables
\end{lstlisting}

The development evidence quoted in this report is committed under \nolinkurl{docs/checkpoints/mg1/runs/}, and
the tables are the CSV files generated by \nolinkurl{make_tables.py}. The held-out data, the evaluator root, model
weights and secrets are not part of the repository.

\chapter{Sample Traces and Screenshots}
\label{app:traces}
\pending{Two or three annotated traces from the development runs are to be added: one run that completes with a
report and no finding, one run in which the agent revises after a refused action, and one run that abstains or is
blocked at preflight. The interface screenshots (timeline, report and review screens) are also to be added. They will
be exported from saved runs so that they match the committed trace files.}

\chapter{Prompts}
\label{app:prompts}
The system prompt and the phase rules of the final prompt version (v9) are reproduced exactly, wrapped for
printing. The JSON schema of the expected reply is appended to the prompt at run time and is not repeated here.

\lstinputlisting[language={},numbers=none,basicstyle=\ttfamily\scriptsize,caption={Final prompt text (version u01\_non\_thinking\_v9)},label=lst:prompt]{chapters/prompt_dump.txt}

\pending{The rubric given to the human scorers is in the repository (\texttt{docs/evaluation/RUBRIC.md}) and is
to be reproduced here once it is frozen with the protocol.}
